\section{Introduction}\label{sec:intro}

The growing global population makes it necessary to develop and improve automatic methods to analyse and recognize activities of people groups in public areas for many reasons including safety (e.g. recent restrictions on meetings due to COVID pandemic) and security (e.g. demonstrations, terrorism, etc.). Usually, human personnel visualize and analyse data from surveillance cameras in order to provide the required actions and decisions depending on the specific problem. This task is very time consuming and not always possible to perform online because, on the one hand, it can be expensive but, above all, the limitations related to human visual inspection: tiredness, fatigue, lack of attention, etc. As a result, camera monitoring is usually consulted after a fact. Recently, however, special attention is being paid to solving these problems with artificial vision and machine learning techniques.

Despite the progress of this area of research, there are many challenges on human activity recognition, mainly related to situations where more than one individual is analyzed. For example, how to distinguish between a street fight and a gather of friends playing, detect an abnormal behavior of an individual among a crowd, or how to analyse a group of people activity over a long period of time. Moreover, the context where the activity takes place has a great impact on the final decision. It allows to determine the type of the activity of if it is normal or abnormal. For instance, a fight in a street is considered abnormal; however, a fight in a boxing ring can be considered as a normal activity. The context defines all the aspects that involve the situation, as the place as we have just mentioned, the time, the conditions, etc. 

Researchers have proposed different strategies to address partial problems of the challenge of analysing group behaviour. In this way, the works propose partial solutions to determine movements, actions, activities of groups, other focus on small groups or crowds, etc. \cite{borja2018short}. Therefore, it is necessary to advance in the proposal of a generic architecture that allows to automate the monitoring of groups of different size and with different levels of semantics.  This motivates our proposal, a multi-stream deep learning architecture able to learn activities of groups using local motion features (Section \ref{sec:parchitecture}). The two main streams of the architecture correspond to the deep-learning variant of the Activity Descriptor Vector (ADV) \cite{Azorin-Lopez2014}. This is a simple descriptor capable of representing the trajectory information of a sequence of images as a collection of local movements that occur in specific regions of the scene. The use of motion features helps to reduce the curse of dimensionality using a deep learning approach. Moreover, 
additional streams are considered to incorporate context information (e.g. location, time, etc.) to enhance the classification of activities. As specific cases, in this paper we instantiate the generic architecture in Section \ref{subsec:arq-mcc} for multi-class analysis with two above-mentioned streams corresponding to the deep-learning ADV variant, with a final decision stage that outperforms state-of-the-art proposals with over $2\%$ accuracy. Also, the architecture is instantiated for one-class classification in Section \ref{subsec:arq-occ}, with three streams including the two ADV variant as in multi-class and a third of context awareness. As it will be in detail presented in Section \ref{sec:exp}, one-class classification instance also improves previous works.

The paper contribution is two-fold:

\begin{itemize}
    \item The main contribution is a multi-stream generic architecture for human group activity recognition, independently of the number of people in the group, that uses local motions features and incorporate context information.  The architecture has been instantiated for multi-class classification and for one-class classification. These two instances include different specific characteristics that, on average, cover a wide range of possibilities.
    
    \item The deep-learning ADV variant (D-ADV) as a simple descriptor that transforms a sequence of color images into a sequence of two-image streams that allow to represent spatio-temporal information using images. They can be used as inputs of deep neural networks models. Compared to end-to-end architectures, our proposal allows to train a classification network from features rather than from raw data reducing the space of solutions and, in consequence, less data is need to train.
\end{itemize}
